%----------------------------------------------------------------------------------------
%       CHAPTER 4
%----------------------------------------------------------------------------------------

\cleardoublepage

\chapterimage{Jay-Matternes---EOCENE.jpg} % Chapter heading image

\chapter{Climates of the Past}\label{ch:paleo-climate}

\hfill \break

\vspace{12mm}

\noindent The Chapter addresses the  climate and ocean circulation states existing at different times during the Phanerozoic and what the primary controls on these state are. (Past ocean biogeochemical cycles will be addressed in a different chapter.) Included are:
\vspace{1mm}

\begin{itemize}
\vspace{1mm}
\item A series of geological time-slices spanning the entire Phanerozoic to provide a primarily tectonic overview of the controls on ocean circulation and climate through time.
\vspace{1mm}
\item A pair of calibrated late Holocene (or pre-industrial) and Last Glacial Maximum (21 ka) configurations.
\vspace{1mm}
\item An early Eocene configuration.
\vspace{1mm}
\item A late Cretaceous configuration (somewhat interchangeable with the early Eocene one in terms of the questions you might pose and hence model experiment that might be conducted).
\end{itemize}

%------------------------------------------------
\newpage
%------------------------------------------------

\section*{Background reading (Phanerozoic climates in general)}
\begin{itemize}
\vspace{1mm}
\item C.R. Scotese, H. Song, B.J. Mills, D.G. van der Meer
Phanerozoic paleotemperatures: the earth’s changing climate during the last 540 million years. Earth Sci. Rev., 215 (2021).
\\\textbf{\href{https://www.sciencedirect.com/science/article/pii/S0024493722002948}{https://www.sciencedirect.com/science/article/pii/S0024493722002948}}
\vspace{1mm}
\item Valdes, P. J., Scotese, C. R., and Lunt, D. J.: Deep ocean temperatures through time, \textit{Clim. Past}, 17, 1483–1506, 2021.
\\\textbf{\href{https://doi.org/10.5194/cp-17-1483-2021}{https://doi.org/10.5194/cp-17-1483-2021}}
\vspace{1mm}
\item Pohl, A., Ridgwell, A., Stockey, R.G. et al. Continental configuration controls ocean oxygenation during the Phanerozoic. \textit{Nature} 608, 523–527 (2022).
\\\textbf{\href{https://doi.org/10.1038/s41586-022-05018-z}{https://doi.org/10.1038/s41586-022-05018-z}}
\end{itemize}

\section*{Background reading (Last Glacial Maximum / deglacial transition)}
\begin{itemize}
\vspace{1mm}
\item Osman, M.B., Tierney, J.E., Zhu, J. et al. Globally resolved surface temperatures since the Last Glacial Maximum. \textit{Nature} 599, 239–244 (2021). \\\textbf{\href{https://doi.org/10.1038/s41586-021-03984-4}{https://doi.org/10.1038/s41586-021-03984-4}}
\vspace{1mm}
\item Osman, M.B., Tierney, J.E., Zhu, J. et al. Globally resolved surface temperatures since the Last Glacial Maximum. \textit{Nature} 599, 239–244 (2021). \\\textbf{\href{https://doi.org/10.1038/s41586-021-03984-4}{https://doi.org/10.1038/s41586-021-03984-4}}
\end{itemize}

\section*{Background reading (Eocene)}
\begin{itemize}
\vspace{1mm}
\item Dunkley Jones, T., D.L. Lunt, D.N. Schmidt, A. Ridgwell, A. Sluijs, P.J. Valdes, and M. Maslin, Climate model and proxy data constraints on ocean warming across the Paleocene-Eocene Thermal Maximum, Earth-Science Reviews 125 123–145 (2013).
\vspace{1mm}
\item Lunt, D. J., Dunkley Jones, T., Heinemann, M., Huber, M., LeGrande, A., Winguth, A., Loptson, C., Marotzke, J., Roberts, C. D., Tindall, J., Valdes, P., and Winguth, C.: A model–data comparison for a multi-model ensemble of early Eocene atmosphere–ocean simulations: EoMIP, Clim. Past, 8, 1717-1736, doi:10.5194/cp-8-1717-2012, 2012.
\vspace{1mm}
\item Tierney JE, Zhu J, Li M, Ridgwell A, Hakim GJ, Poulsen CJ, Whiteford RDM, Rae JWB, Kump LR. Spatial patterns of climate change across the Paleocene-Eocene Thermal Maximum. \textit{Proc Natl Acad Sci}, 2022. 
\\\textbf{doi:10.1073/pnas.2205326119}
\end{itemize}

\section*{Background reading (Cretaceous)}
\begin{itemize}
\vspace{1mm}
\item Barron, E.: A warm equable Cretaceous: the nature of the problem, Earth-Science Reviews 19, 305–338, 1983.
\vspace{1mm}
\item Huber, B. T., Norris, R. D., and MacLeod, K. G.: Deep-sea paleotemperature record of extreme warmth during the Cretaceous, Geology 30, 123–126, 2002.
\vspace{1mm}
\item Donnadieu, Y., et al., Modelling the primary control of paleogeography on Cretaceous climate, Earth and Planetary Science Letters 248, 426–437, 2006.
\vspace{1mm}
\item Hunter, S. J., et al., Modelling Maastrichtian climate: investigating the role of geography, atmospheric \(CO_{2}\) and vegetation, Clim. Past Discuss. 4, 981–1019, 2008.
\vspace{1mm}
\item Jenkyns, H. C., Forster, A., Schouten, S., and Sinninghe Damste, J. S.: High temperatures in the Late Cretaceous Arctic Ocean, Nature 432, 888–892, 2004.
\end{itemize}

%------------------------------------------------
\newpage
%------------------------------------------------

\section*{READ.ME}

You will need to download new \textit{restart} files for this Chapter (both 10,000 year spin-ups) depending on what time intervals you are interested in and hence experiments you want to conduct.

\vspace{2mm}

The \textit{re-start} files for the Phanerozoic series are of the form:

\begin{enumerate}[noitemsep]
\vspace{2mm}
\item[\textbf{[time yyy]:}]
\footnotesize\begin{verbatim}
$ wget --no-check-certificate http://www.seao2.info/cgenie_output/muffin.C.yyyrdP1_.r.SPIN.tar.gz
\end{verbatim}\small
\end{enumerate}
\vspace{2mm}

where \texttt{yyy} (in the string: \texttt{yyyrdP1\_.r}) is a code for the time slice, in Ma. The options are:

\vspace{1mm}
\small\begin{itemize}[noitemsep]
\item []
\begin{itemize}[noitemsep]
\item \texttt{0\_\_} == 0 Ma
\item \texttt{20\_} == 0 Ma
\item [] ...
\item \texttt{100} == 100 Ma
\item \texttt{120} == 120 Ma
\item [] ...
\item \texttt{520} == 520 Ma
\item \texttt{540} == 540 Ma
\end{itemize}
\end{itemize}\normalsize

\vspace{2mm}

For the more focussed time intervals and questions: 

\begin{enumerate}[noitemsep]
\vspace{2mm}
\item[\textbf{Holocene:}]
\footnotesize\begin{verbatim}
$ wget --no-check-certificate http://www.seao2.info/cgenie_output/muffin.C.GIteiiaa.r.SPIN.tar.gz
\end{verbatim}\small
\vspace{2mm}
\item[\textbf{LGM:}]
\footnotesize\begin{verbatim}
$ wget --no-check-certificate http://www.seao2.info/cgenie_output/muffin.C.GIteiiva.r.SPIN.tar.gz
\end{verbatim}\small
\vspace{2mm}
\item[\textbf{Eocene:}]
\footnotesize\begin{verbatim}
$ wget --no-check-certificate http://www.seao2.info/cgenie_output/LAB_5b.SPIN.tar.gz
\end{verbatim}\small
\vspace{2mm}
\item[\textbf{Cretaceous:}]
\footnotesize\begin{verbatim}
$ wget --no-check-certificate http://www.seao2.info/cgenie_output/LAB_5.SPIN.tar.gz
\end{verbatim}\small
\end{enumerate}
\vspace{2mm}

Extract the results in the usual way and in the usual place ... and return to \texttt{genie-main}.

%------------------------------------------------
\newpage
%------------------------------------------------

\section{Ocean circulation through time}

We start this Chapter from the very (almost) broadest perspective and consider how and why climate, and particularly global ocean circulation, has changed through the Phanerozoic.

The \textbf{muffin} configurations used here are all based on \textit{Pohl et al.} [2022], which looked at the role of tectonics, and specifically the position of the continental plates, had on patterns on ocean circulation and the consequences of this for marine biogeochemical cycles and specifically oxygenation of the oceans. However, these have been simplified (the full configurations can be found via the model code availability information in the paper itself) as climate-only configurations (for speed), but including an age tracer.

The specific series of configurations from \textit{Pohl et al.} [2022] are those of 'series \#1' (see paper) in which the solar constant and atmospheric \(CO_{2}\) (here, converting to radiative forcing) are invariant across the entire series of time-slices. (The purpose of this was to expose the role of tectonics-only in modulating global climate and ocean circulation patterns and hence adopted the same short and long-wave forcing boundary conditions.)

A single template \textit{user-config} is provided in the \textsf{\footnotesize EXAMPLES} sub-directory of \textsf{\footnotesize genie-userconfigs}: \textsf{\footnotesize muffin.C.yyyrdP1\_.r.SPIN} and which can be used to create an experiment for any of the different Phanerozoic time-slices. In contrast, there a base-config for each of the time-slices (becasue different model boundxary conditions need to be specified for each). These adopt the form: \textsf{\footnotesize muffin.C.yyyrdP1\_.r.config} where \textsf{\footnotesize yyy} is the geological time (in Ma) code: 

\vspace{1mm}
\small\begin{itemize}[noitemsep]
\item []
\begin{itemize}[noitemsep]
\item \texttt{0\_\_} == 0 Ma
\item \texttt{20\_} == 0 Ma
\item [] ...
\item \texttt{100} == 100 Ma
\item \texttt{120} == 120 Ma
\item [] ...
\item \texttt{520} == 520 Ma
\item \texttt{540} == 540 Ma
\end{itemize}
\end{itemize}\normalsize

\vspace{2mm}
Experiments can then be run as follows (taking the 0Ma time-slice as an example):

\vspace{-2mm}\small
\begin{verbatim}
$ ./runmuffin.sh muffin.C.0__rdP1_.r EXAMPLES muffin.C.yyyrdP1_.r.SPIN 10000
\end{verbatim}
\normalsize\vspace{-2mm}

\noindent although you might want to make use of the re-starts and only run for much shorter, e.g.

\vspace{-2mm}\small
\begin{verbatim}
$ ./runmuffin.sh muffin.C.0__rdP1_.r EXAMPLES muffin.C.yyyrdP1_.r.SPIN 100
        muffin.C.0__rdP1_.r.SPIN
\end{verbatim}
\normalsize\vspace{-2mm}

\noindent ... and of course creating your own \uline{distinct name} (and copied/re-named \textit{user-config}) for the experiment (which you might also move to a different directory), e.g.

\vspace{-2mm}\small
\begin{verbatim}
$ ./runmuffin.sh muffin.C.0__rdP1_.r LABS experiment_0Ma 100
        muffin.C.0__rdP1_.r.SPIN
\end{verbatim}
\normalsize\vspace{-2mm}

\noindent Make sure that the time code in the \textit{base-config} and \textit{re-start} match!

\noindent But also, make sure that your chosen \textit{user-config} name is not the same as the \textit{re-start} name ... (In the example, \textsf{\footnotesize muffin.C.yyyrdP1\_.r.SPIN} is renamed to \textsf{\footnotesize experiment\_0Ma}, and then uses the restart: \textsf{\footnotesize muffin.C.0\_\_rdP1\_.r.SPIN}.)

\vspace{1mm}
\noindent\rule{4cm}{0.5pt}
\vspace{2mm}

\noindent So ... what might you actually 'do' with all these configurations (refer to e.g. \textit{Pohl et al.} [2022] and \textit{Valdes et al.} [2021] for ideas of how series of model configurations and experiment time-slices have been analysed and presented before)? The following sub-sections provide some ideas:

%------------------------------------------------
\newpage
%------------------------------------------------
%
\subsubsection{Questioning configuration fidelity (resolution and re-gridding)}

\vspace{2mm}

The \textbf{muffin} configurations are re-gridded from a relatively low resolution coupled GCM (for details, see: \textit{Pohl et al.} [2022]) which itself takes the continental configuration from a much higher resolution plate reconstruction. The first thing to note then is that re-gridding form relatively low, to even lower, may result in key features of the modern world, e.g. ocean gateways, becoming poorly resolved.

So to start, why not create a new experiment (and \textit{user-config} file) run on the 0 Ma time-slice for 100 years. Analyse patterns of SST, baroclinic circulation (and perhaps also plot some velocity vectors) to assess  something of the strength and position of the ACC,  barotropic (global overturning) circulation (Atlantic and Pacific basin MOCs are not provided as output), and the ventilation age at depth (either  the 2D netCDF benthic layer output or taking a deep slice from the 3D data field). You can contrast the outputs, either visually, or perhaps creating a difference map, with the standard modern configuration (\textsf{\footnotesize p\_worjh2} analysed in the earlier, ocean circulation Chapter -- \textit{re-start}: \textsf{\footnotesize muffin.C.p\_worjh2.r.SPIN}).

However, poorly reproduced features of the modern global ocean circulation and climate may also arise because in the common climate forcing assumptions on the Phanerozoic series of experiments, atmospheric \(CO_{2}\) was assumed to be approximately \(\times8\) modern. An obvious next experiment to run is then to create an experiment (and new \textit{user-config}) with the radiative forcing scaling set:

\vspace{-2mm}\footnotesize
\begin{verbatim}
# scaling for atmospheric CO2 radiative forcing, calculated as the multiplier of 278 ppm CO2
ea_radfor_scl_co2=1.0
\end{verbatim}
\normalsize\vspace{-2mm}

The re-start is of a warm climate state, and although the ocean should relatively quickly cool (compared to the longer time-scale of top-down warming), you will need to run your experiment for rather longer. How long? Something to test and explore for yourself, perhaps by running an experiment longer than you think you'll need, and analysing how quickly different components of the climate system and ocean depth, approach an equilibrium state. (Having obtained this knowledge once, you can judge a sufficient run-time and use this in any similar subsequent experiment.)

When this cooler climate state is attained, you can repeat the previous analysis. When you do this, you may find 2 continuing deficiencies in the simulation of ocean circulation:

\begin{enumerate}[noitemsep]
\vspace{1mm}
\item A non-existent AMOC.
\vspace{1mm}
\\Q. How can you tell if you do not have access to an AMOC netCDF output?
\\A. Analyse the distribution of ventilation age through the global deep sea -- compare with the standard modern configuration as well as using your understanding of modern ocean circulation patterns and what this should mean for where the youngest vs. oldest deep water should be found.
\\How to fix this?
\\If you were doing thing 'properly', you could create a 'mask' to define a region where freshwater water will be removed from the surface in one region (the North Atlantic) and added elsewhere (e.g. the North Pacific). This is very do-able, but a little too involved at this stage.
\\How about a anti-hosing \textit{a-la} the experiments in the earlier ocean circulation Chapter? i.e. applying a positive salinity flux at some promising location i the North Atlantic?

\newpage

\noindent To set this up, simply add the parameter requesting a flux \textit{forcing} to your \textit{user-config}:

\vspace{-0mm}\footnotesize
\begin{verbatim}
bg_ctrl_force_GOLDSTEInTS=.true.
bg_par_forcing_name="pyyyyz.Fsal"
bg_par_force_point_i=22
bg_par_force_point_j=33
bg_par_force_point_k=16
bg_par_ocn_force_scale_val_2=-2.0E17
\end{verbatim}
\normalsize\vspace{-1mm}

but remembering that a negative scaling value is giving you negative salinity (i.e. fresh water), and you want salt instead (i.e. positive). Also note that the model grid is different.\footnote{The distribution of land vs. ocean and the depth of the ocean is defined by a simple text file with the extension \textsf{\footnotesize .k1}. For the 0 Ma configuration, can find this file in: \textsf{\footnotesize genie-paleo/0\_\_rdP1\_}. Simply count columns to the right (ignoring the first one) for the \texttt{i} value, and rows up (ignoring the first one) for the \texttt{j} value.}

\vspace{1mm}
\item A weak ACC.
\vspace{1mm}
\\The ACC in the standard model is relatively weak (you can find its strength, in units of Sv, from the 2D netCDF baroclinic output). (The 'real' ACC is something like \(120Sv\), although you might Google and check this value. Even read a paper!!!)
\\If you find that the ACC in your cool climate, \textsf{\footnotesize 0\_\_rdP1\_} experiment is much weaker ... what to do about it?
\\In the model, the principal controls on the strength of the ACC are:
\vspace{1mm}
\begin{enumerate}[noitemsep]
\item The high southern latitude wind stress.
\vspace{1mm}
\\Wind stress fields can be adjusted, but the format and the fact that the wind stress boundary conditions are provided on 2 different model grids, makes adjusting specific the values in restricted regions or latitude bands, tiresome. So you could try adjusting the model parameter that scales the values in the boundary condition files, to what is actually applied in the model at the ocean surface. For equally tiresome reasons, this scaling value is duplicated across 2 different parameters, and the value of both must be changed together.
\vspace{1mm}
\\The parameters (and default values) are:
\vspace{-1mm}\footnotesize
\begin{verbatim}
# scl_tau / SclTau
go_13=1.531013488769531300
# scl_tau (should be same as GOLDSTEIN's value) / SclTau
ea_11=1.531013488769531300
\end{verbatim}
\normalsize\vspace{-1mm}
These parameter value definition lines normally live in the \textit{base-config} file, but you can simply copy them into your \textit{user-config} and increase the values. (by how much? something try and find out!)
\\Don't forget to run a control experiment in which you also copy these parameter lines into a (second) \textit{user-config} file ... but \uline{don't change their values}.

\vspace{1mm}
\item The width and sill depth of key gateways.
\vspace{1mm}
\\The Drake Passage is likely to be the critical gateway restricting the ACC if you find it far too weak. If you remember, the distribution of land vs. ocean and the depth of the ocean (as well as the direction of run-off on land), is defined by a simple text file with the extension \textsf{\footnotesize .k1}. For the 0 Ma configuration, can find this file in: \textsf{\footnotesize genie-paleo/0\_\_rdP1\_}. \textbf{Make a backup copy before you change anything!} \\Unlike in creating new \textit{user-config} files, you need to edit the original file (and filename).\footnote{Otherwise you have to change the corresponding filenames in the \textit{base-config}.} In editing the file, remember that lower integer values correspond to the ocean being deeper at that location.
\vspace{1mm}
\begin{enumerate}[noitemsep]
\item First try a simple deepening -- i.e. you are not changing where land vs. ocean is, simply how deep the ocean is.
\\Don't forget a control experiment ...
\item Then, with some \uline{care}, you could try widening the gateway. Land is a number in the \(90s\), and ocean is a number \(\leqslant16\). However, there is a second file involved that requires consistent editing, which has a hard-to-decipher format.
\\To widen the gateway, you are going to have to use \textbf{muffingen}, use the \textsf{\footnotesize .k1} file as an input, and then replace both the current \textsf{\footnotesize .k1} and \textsf{\footnotesize .psiles} files.
\end{enumerate}
 
\end{enumerate}
 
\end{enumerate}

\vspace{2mm}
After all that ... did you manage to generate an 'improved' modern global circulation at all?

%------------------------------------------------
\vspace{1mm}
\noindent\rule{4cm}{0.5pt}
%
\subsubsection{Role of continental configuration (and sealevel) in modulating global climate}

\vspace{2mm}

You could simply take the re-starts, or create new experiments and \textit{user-configs} and run selection (or all) of the different configurations with their parameter values unchanged. Try and generate a similar (or better!) analysis to \textit{Pohl et al.} [2022], of how key climate variables change with time through the Phanerozoic.

%------------------------------------------------
\vspace{1mm}
\noindent\rule{4cm}{0.5pt}
%
\subsubsection{Role of evolving solar constant in modulating global climate}

\vspace{2mm}

In the study of \textit{Pohl et al.} [2022], the solar constant was deliberately kept at its modern value. However, in reality it would have increased during the Phanerozoic and at 540 Ma would have been about \(1306Wm^{-2}\) compared with \(1368Wm^{-2}\) today. The parameter controlling this is:

\vspace{-2mm}\footnotesize
\begin{verbatim}
# Solar constant (W m-2)
ma_genie_solar_constant=1368.0
\end{verbatim}
\normalsize\vspace{-2mm}

\noindent(with its default, modern value). 
You know what to change the value to for 540 Ma. What about the geological time in between? You could make a linear change with time assumption. You could also dig out the literature and calculate a more 'exact' value -- see \textit{Feulner} [2012]\footnote{Feulner, G. (2012), The faint young Sun problem, Rev. Geophys., 50, RG2006, \textbf{doi:10.1029/2011RG000375}.}. In fact you might try both and see whether a linear assumption is a good approximation (or not) compared to the nonlinear function of \textit{Feulner} [2012].

%------------------------------------------------
\vspace{1mm}
\noindent\rule{4cm}{0.5pt}
%
\subsubsection{Role of evolving CO2 with time}

\vspace{2mm}

Atmospheric \(pCO2\) also changes with time. Again, \textit{Pohl et al.} [2022] fixed with (at a value equivalent to \(x8\)).

So in addition to having pre-industrial radiative forcing (\(x1\)) set for for 0 Ma, you might create some new experiments for some (or all!) of the older time-slices with more appropriate \(CO_{2}\) (or rather: radiative forcing) values. And where are you going to get earlier-in-geological time \(CO_{2}\) values from ... ? Some possibilities (and including tracing down the data-sets provided in association with the publications, although you might simply estimate values from published curves and individual points in publication figures):

%------------------------------------------------
\newpage
%------------------------------------------------

\begin{itemize}[noitemsep]
\vspace{1mm}
\item Foster, G., Royer, D. \& Lunt, D. Future climate forcing potentially without precedent in the last 420 million years. \textit{Nat Commun} 8, 14845 (2017). \\\textbf{\href{https://doi.org/10.1038/ncomms14845}{https://doi.org/10.1038/ncomms14845}}
\vspace{1mm}
\item D.L. Royer. CO2-forced climate thresholds during the Phanerozoic.
\textit{Geochimica et Cosmochimica Acta}, 70 (2006), pp. 5665-5675.
\end{itemize}

\vspace{1mm}
Also, see: \href{https://www.paleo-co2.org/}{\textbf{https://www.paleo-co2.org/}}

\vspace{1mm}
You are then going to have to convert the \(CO_{2}\) value, to a multiplier compared to \(278 ppm\) in order to deduce the scaling value needed in the \textit{user-config}.

\vspace{1mm}
How does this now affect the climate and ocean circulation state of a specific point in geological time, or affect trends through the Phanerozoic?

%------------------------------------------------
\vspace{1mm}
\noindent\rule{4cm}{0.5pt}
%
\subsubsection{Evolving solar constant \uline{and} CO2!}

\vspace{2mm}

Obviously, both forcings on the energy balance of the climate system change with time through the Phanerozoic, so you might want to repeat experiments and change he value of both. 

%------------------------------------------------
\vspace{1mm}
\noindent\rule{4cm}{0.5pt}
%
\subsubsection{Circulation tipping points}

\vspace{2mm}

As per \textit{Pohl et al.} [2022], some of the deeper time (Paleozoic) configurations give rise to stable oscillations in ocean circulation. These states are not only a function of tectonic arrangement, but climate. The question is simply then whether you can find climate states -- different values of solar constant and/or radiative forcing (as compared to the defaults) -- that make the oscillations 'go away'. Equally, you can find different climate states that induce oscillations for any of the continental configuration that by default do not exhibit oscillations?

\vspace{1mm}
Note that these experiment might take some time ... typically, a minimum of about 10,000 model years needs to be run in order to identify the presence/absence of an oscillatory state.

%------------------------------------------------
\newpage
%------------------------------------------------

\section{The LGM and last deglaciation}

Two example configurations and \textit{re-starts} are given -- one for the late Holocene (effectively, modern (or preindustrial)), and one for the last Glacial Maximum (LGM) (nominally at \(21ka\)). Both model configurations are derived using \textbf{muffingen} from fully coupled GCM experiments created with the HadCM3 fully coupled GCM, and hence \textbf{muffin} is forced by 2D stress and wind speed, plus a zonal mean albedo profile, all originally simulated by a GCM.

\vspace{1mm}
The two configurations\footnote{Note that they differ only minutely in their filename ...} are:

\begin{enumerate}[noitemsep]
\vspace{1mm}
\item \texttt{muffin.C.GIteiiaa.r.SPIN}
\vspace{1mm}
\\This experiment is additionally configured with radiative forcing corresponding to late Holocene atmospheric \(CO_{2}\) and and modern orbits.
\\A plausible circulation in the Atlantic (and good statistical fit to available benthic $\delta^{13}$C data) is achieved with a modern-like flux correction in \textbf{muffin}:
\vspace{-2mm}\begin{verbatim}
bg_par_ocn_force_scale_val_2=0.30
\end{verbatim}\vspace{-2mm}
i.e. a flux of \(0.30 Sv\).
\vspace{1mm}
\item \texttt{muffin.C.GIteiiva.r.SPIN}
\vspace{1mm}
\\This experiment is additionally configured with radiative forcing corresponding to LGM atmospheric \(CO_{2}\) and has a 21 ka orbital configuration. Mean ocean salinity is increased by 3\% to reflect the decrease in ocean volume due to the removal of freshwater to the ice sheets.
\\A plausible circulation in the Atlantic (and good statistical fit to available LGM benthic $\delta^{13}$C data) is achieved with a reduced flux correction in \textbf{muffin}:
\vspace{-2mm}\begin{verbatim}
bg_par_ocn_force_scale_val_2=0.05
\end{verbatim}\vspace{-2mm}
i.e. a flux of \(0.05 Sv\).
\end{enumerate}
\vspace{1mm}

Both configurations include a number 'age' tracer that records the mean time since a parcel of water last 'saw' the ocean surface -- the ventilation age.

\vspace{1mm}
Note that the \textit{user-configs} are not in the sub-directory \textsf{\footnotesize LABS}, but instead in \textsf{\footnotesize EXAMPLES}. You can either copy them to \textsf{\footnotesize LABS} and run as normal, or run them from \textsf{\footnotesize EXAMPLES} by changing the \textit{user-config} sub-directory you specify at the command line when running an experiment (better to copy). Make sure that you change the name of the \textit{user-configs} you run so that they are not the same as the \textit{re-start} experiment names ...

\vspace{1mm}
You can run the Holocene and LGM experiments on from their corresponding \textit{re-start} (and assuming that you have renamed the \textit{user-configs} to \textsf{\footnotesize experiment\_HOL} and \textsf{\footnotesize experiment\_LGM}, respectively)\footnote{Make sure that you copy and re-name the \textit{user-config} files otherwise you'll over-write the re-starts that you just downloaded!} by e.g.:

\small\vspace{-2pt}\begin{verbatim}
$ ./runmuffin.sh muffin.C.GIteiiaa.r EXAMPLE experiment_HOL 100 
        muffin.C.GIteiiaa.r.SPIN 
$ ./runmuffin.sh muffin.C.GIteiiva.r EXAMPLES experiment_LGM 100 
        muffin.C.GIteiiva.r.SPIN
\end{verbatim}\vspace{-2pt}\normalsize

%------------------------------------------------
\newpage
%------------------------------------------------

Experiments are a little more challenging to really get anything out of because the model steady-state take ca. 5000 or so years to come close to steady-state. Hence it is unlikely that you will see the full impact of any parameter change unless you run an experiment overnight on the queue. 50-100 years would probably the minimum to start to see meaningful adjustments to e.g. ocean circulation.

\vspace{1mm}
\noindent\rule{4cm}{0.5pt}
\vspace{2mm}

\noindent There are fewer obvious interesting experiments for the late Holocene configuration and re-start, but you might still explore some perturbations of ocean circulation on the basis that most (almost without exception all) model de-glacial/AMOC collapse experiments are run and published based on a modern or slightly modified modern configuration. SO for instance, you might set the flux adjustment that maintains the AMOC to zero:
\vspace{-2mm}\begin{verbatim}
bg_par_ocn_force_scale_val_2=0.0
\end{verbatim}\vspace{-2mm}
and see what happens. This should look pretty much like the freshwater hosing experiments (i.e. rather than apply fresh water to the North Atlantic region in a model that already has a constant prescribed net salinity transport from North Pacific to North Atlantic, you are simply taking away the underlying prescribed transport).

\vspace{1mm}
For the LGM configuration, you might also try completely collapsing the AMOC (reducing the flux correction from \(0.05 Sv\) to zero (or even slightly negative). You are thereby simulating an AMOC collapse from a glacial state -- aka 'H0'. Note that as an alternative to simply removing the default flux adjustment, it is possible to create a time-dependent flux of freshwater (as traditionally applied in models, e.g. a pulse or ramp up/down) by editing:
\vspace{-2mm}\begin{verbatim}
biogem_force_flux_ocn_sal_sig.dat
\end{verbatim}\vspace{-2mm}
in the \textit{forcing}.

\vspace{1mm}
You might also consider what a slightly strong LGM AMOC might look like in terms of the pattern of benthic ventilation ages. A moderate and reasonable increase in the flux adjustment would be to dial it up to \(0.10 Sv\). (Dialling all the way up to \(0.30 Sv\) would enable you to access what a modern-like AMOC would look like in an otherwise glacial world.)

%------------------------------------------------
\newpage
%------------------------------------------------

\section{The climate of the early Eocene}

A previously spun-up state of early Eocene   (ca. 56 Ma) climate (\textsf{\footnotesize LAB\_5b.SPIN}) plus example \textit{user-config} (\textsf{\footnotesize LAB\_5b.EXAMPLE}) are provided, and run e.g.:
\vspace{-1mm}\small\begin{verbatim}
$ ./runmuffin.sh cgenie.eb_go_gs_ac_bg.p0055c.NONE LABS
 LAB_5b.EXAMPLE 10 LAB_5b.SPIN
\end{verbatim}\normalsize\vspace{-1mm}
(Note that for speed, no carbon cycle is configured.)

A second configuration is provided that includes the ocean dye (color) tracer. Using the same \textit{re-start}, the corresponding \textit{user-config} (\textsf{\footnotesize LAB\_5b.colorinjection}) for this can be run:
\vspace{-1mm}\small\begin{verbatim}
$ ./runmuffin.sh cgenie.eb_go_gs_ac_bg.p0055c.rb LABS
 LAB_5b.colorinjection 10 LAB_5b.SPIN
\end{verbatim}\normalsize\vspace{-1mm}

\vspace{1mm}
\noindent\rule{4cm}{0.1mm}
\vspace{2mm}

\noindent What to ‘do’ with this? Again – start by tracing and analyzing the patterns of ocean circulation in the Eocene world, as per how you did for the modern world. Where does deep water form in the Eocene? Is the surface ocean circulation as expected?
Can you see any patterns emerging? At least in model world – what seems to be dictating the locations of deep water formation? Does overturning and/or surface circulation appear to significantly vary as a function of continental configuration or radiative forcing? If not, why not?

\vspace{1mm}
\noindent\rule{4cm}{0.1mm}
\vspace{2mm}

\noindent But how do you know how close, or not, the model climate is to ‘reality’?
Available from \href{https://www.seao2.info/mymuffin.html}{my website}, on the LH side of the page under ‘got data?’ and at the bottom of the list, are 2 netCDF files containing proxy sea surface temperature (SST) data re-gridded to the cGENIE model grid. One is for pre-PETM SSTs, and one for peak PETM SSTs. (The data are from \textit{Dunkley Jones et al.} [2013].) 
These netCDF files can be opened and visualized in exactly the same way in exactly the same way as per the cGENIE output files. You’ll see a few (there are not many!) data points, in a sea of grey (which stands for ‘no value’). (Note that it is best to turn data interpolation OFF.)

More useful … is that the data distribution can be combined with the model ocean temperature output field, in a difference map (see earlier for notes on creating difference maps). You can thus visualize the pattern of model-data mismatch. (Note you might want to hit ‘Fit to data’ to autoscale the plot, or simply pick and manually enter a +/- limit for the data plotting, perhaps ensuring it is symmetrical about zero such that red anomalies could represent proxy SST values higher than in the model, and blue point data being lower than in the model.)

Furthermore, in the Arrays tab, you can switch to a Zonal Average view rather than a Map view (...it is probably easier to visualize the model-data misfit in this way.)

\vspace{1mm}
\noindent\rule{4cm}{0.1mm}
\vspace{2mm}

\noindent So, starting with the pre-PETM data and given model configuration, make some assessment of how the model fits the data (or not). What might be the reason for the misfit? You might test adding, and adjusting, some of the parameter values controlling surface SST from the Cretaceous climate exercise (next sub-section). Can you reduce the model-data misfit?

\vspace{1mm}
\noindent\rule{4cm}{0.1mm}
\vspace{2mm}

\noindent Finally, the PETM SST data given in the 2nd netCDF file are from the peak of the PETM. The ocean is rather warmer as compared to just prior to the PETM. Your task is to determine (either using the default model configuration, or your adjusted configuration) much radiative forcing (parameter \texttt{ea\_radfor\_scl\_co2}) and hence \(CO_{2}\) is required to explain this warming?

%------------------------------------------------
\newpage
%------------------------------------------------

\section{The climate of the Cretaceous}

A previously spun-up state of Maastrichtian climate with \textbf{muffin} in a ca. 70 Ma configuration\footnote{Note that for speed, no carbon cycle is selected in this configuration.} (\texttt{LAB\_5.SPIN}) is provided as a starting point. A \textit{user-config} (\texttt{LAB\_5.EXAMPLE}) that continues on this climate state is run:

\vspace{-2mm}\small\begin{verbatim}
$ ./runmuffin.sh cgenie.eb_go_gs_ac_bg.p0067f.NONE LABS
 LAB_5.EXAMPLE 10 LAB_5.SPIN
\end{verbatim}\normalsize\vspace{-2mm}

Your task now is ... ‘simple’: Account for the Cretaceous reduced Equator-to-pole surface temperature gradient (compared to modern), particularly the apparently much warmer poles. Different data-based time-slices (including the Maastrichtian) are provided in \textit{Huber et al.} [2002], although the low latitude \(\delta^{18}O\) based temperatures are now not considered reliable. \textit{Jenkyns et al.} [2004] contains high latitude (Arctic) data for the Maastrichtian. There are proxy-derived latitudinal temperature gradients and model-data studies for earlier in the Cretaceous – the problem is essentially the same.

The 2-D NetCDF results file contains the surface air temperature field (and sea-ice cover, if any). The 3-D NetCDF results file contains fields for ocean temperatures (and salinity). Both contain continental configuration and ocean bathymetry.

\textbf{Panoply} will plot the zonal average for you (as used in model-data comparisons – e.g., see Bice and Norris [2002]) – in the Array(s) tab, the Plot can be set to Zonal Averages rather than Map. You can get a smooth curve by selecting Interpolate. Remember you can set (and fix) scales rather than let Panoply auto-scale continually.

\vspace{1mm}
\noindent\rule{4cm}{0.1mm}
\vspace{2mm}

\noindent The following ‘controls’ over the climate system are provided to you in the form of a list of parameters at the bottom of the \texttt{LAB\_5.EXAMPLE} \textit{user-config} file for editing:

\vspace{-1mm}\footnotesize\begin{verbatim}
# --- MISC -----------------------------------------------
#
# === ATMOSPHERE ===
# CO_2 radiative forcing scaling factor [DEFAULT = 4.0]
ea_radfor_scl_co2=4.0
# CH4 radiative forcing scaling factor [DEFAULT = 1.0]
ea_radfor_scl_ch4=1.0
# Equator-to-pole different in planetary albedo [DEFAULT = 0.260]
ea_albedop_amp=0.260
# Baseline planetary albedo [DEFAULT = 0.200]
ea_albedop_offs=0.200
# atmospheric diffusivity of temperature (horizontal) [DEFAULT = 5.0e6]
ea_12=5.0e6
# === OCEAN ===
# ocean diffusivity of temperature + salinity (horizontal) [DEFAULT = 1494.4]
go_14=1494.4
# scaling for wind stress (set values of both identical) [DEFAULT = 1.531]
go_13=1.531
ea_11=1.531
\end{verbatim}\normalsize\vspace{-1mm}

%------------------------------------------------
\newpage 
%------------------------------------------------

Most of these parameters are associated with the radiative forcing of climate or atmospheric transports. The most useful ones are likely to be:

\begin{itemize}[noitemsep]

\vspace{1mm}
\item[\textbf{(i)}] The line:
\vspace{-1mm}\begin{verbatim}
ea_radfor_scl_co2=4.0
\end{verbatim}\vspace{-1mm}
which specifies a radiative forcing of climate by \(CO_{2}\) equivalent to \(4\) times modern \(CO_{2}\) (i.e., \(4\times278 ppm = 1112 ppm\)) as per before (e.g. snowball Earth experiments).
\\The line:
\vspace{-1mm}\begin{verbatim}
ea_radfor_scl_ch4=1.0
\end{verbatim}\vspace{-1mm}
specifies a radiative forcing of climate by \(CH_{4}\) equivalent to \(1\) times modern \(CH_{4}\) (i.e., \(1\times700 ppb\)).\footnote{But effectively, this is going to do exactly the same as changing radiative forcing due to \(CO_{2}\).}

\vspace{1mm}
\item[\textbf{(ii)}] cGENIE, as configured here, does not have a land surface scheme (no snow cover) nor clouds nor ice sheets, so a planetary albedo is prescribed (see Figure 4.1). This varies with latitude and is parameterized after a fully coupled GCM simulation. There is a parameter which controls how the albedo varies as a function of latitude, which can be adjusted:
\vspace{-1mm}\begin{verbatim}
ea_albedop_amp=0.260
\end{verbatim}\vspace{-1mm}
However, if you vary this, why are you doing it? (What is the physical justification for giving the poles a higher or lower albedo relative to the Equator?)
There is also a parameter that sets the baseline (minimum) albedo:
\vspace{-1mm}\begin{verbatim}
ea_albedop_offs=0.200
\end{verbatim}\vspace{-1mm}
i.e., albedo is primarily a sum of the baseline value plus the Equator-to-pole slope (times the latitude).

\vspace{1mm}
\item[\textbf{(iii)}] The diffusivity of heat in the atmosphere (since it is a simple 2-D atmospheric model, with no atmospheric circulation, the atmosphere is made ‘diffusive’ to help capture heat and moisture transport) is:
\vspace{-1mm}\begin{verbatim}
ea_12=5.0e6
\end{verbatim}\vspace{-1mm}
Note the maximum value the model can cope with is ca. \(1.0E7\).

\end{itemize}

\begin{figure}
\includegraphics[width=0.6\textwidth]{tmp-albedo.png}\centering
\vspace{-6mm}
\caption{Prescribed planetary albedo.
The latitudinal (from 90$^{\circ}$S (-90$^{\circ}$N) on the left, to 90$^{\circ}$N on the right) profile of planetary albedo as calculated in a fully coupled GCM is given in green, and the cGENIE ‘fit’ in black.}
\label{fig:tmp-albedo}
\end{figure}

\vspace{1mm}
\noindent\rule{4cm}{0.1mm}
\vspace{2mm}

\noindent Run the model for however long you think is ‘necessary’ (/justified). The surface climate will approach equilibrium ‘relatively’ quickly. Deep ocean temperatures will typically take thousands of years to fully adjust ... You can assess how the model approaches equilibrium most easily from the atmospheric temperature time-series results file, and from the ocean temperature time-series results file (allowing to you to contrast surface and whole ocean temperature changes).

\vspace{1mm}
\noindent\rule{4cm}{0.1mm}
\vspace{2mm}

\noindent You can also try tracing and analyzing the patterns of ocean circulation in the Cretaceous world, in the same way as you did for the modern climate system. And in fact, it would be a useful exercise to directly compare your previous modern results vs. Cretaceous results, or even better, carry out a set of paired experiments (e.g., similar locations and fluxes of dye injection) for both modern and Cretaceous. (Note that if you do this and swap back and forth between modern and Cretaceous configurations, you cannot submit straight away to the cluster but must first briefly run the experiment at the command line in order that cGENIE is compiled into the new continental configuration.)

A new \textit{user-config} (\texttt{LAB\_5.colorinjection}) that continues on from this climate state, can be run:
\vspace{-1mm}\begin{verbatim}
$ ./runmuffin.sh cgenie.eb_go_gs_ac_bg.p0067f.rb LABS
 LAB_5.colorinjection 10 LAB_5.SPIN
\end{verbatim}\vspace{-1mm}
(Also note that the \textit{base-config} is different as it now includes the definition of 2 dye tracers.)

The default locations for the dye release are set different compared to in the modern configuration. BUT, the default location is not necessarily ideal / particularly revealing … so you’ll need to look count grid cells West-to-East (the i direction) and from South to North (the j direction) in order to determine a more suitable location for tracing circulation.

By means of dye tracing, looking at the (global only) overturning stream-function, and/or temperature and salinity (and density) profiles, see if you can identify where in the Cretaceous ocean deep water forms in the model. AND, more importantly, think about WHY does it form where it does?

Also: what about surface ocean circulation and gyres? Are these located where you would expect (e.g. based on your understanding of modern ocean circulation). Plot velocity vectors to help, and/or refer to the Barotropic streamfunction output in the 2D netCDF file.

%----------------------------------------------------------------------------------------
%----------------------------------------------------------------------------------------